
% IMPORTANTE:
% Este contenido procede de una política docente reutilizable,
% pero debe revisarse según la institución, el programa y las
% características particulares de cada curso.

% ============================================================
% OTROS ASPECTOS DEL CURSO
%
% Esta sección reúne las reglas de asistencia, comunicación,
% honestidad académica y uso de inteligencia artificial.
% ============================================================

\section{Otros aspectos para tener en cuenta}


% ------------------------------------------------------------
% ASISTENCIA
% ------------------------------------------------------------

% \subsection crea un nivel inferior a \section.
\subsection{Asistencia}

El curso requiere un mínimo del 85\,\% de asistencia a las sesiones
hábiles programadas. Este cálculo excluye festivos, jornadas de la
Semana Universitaria y sesiones cuya no asistencia sea acordada
formalmente con el docente.

Se considerará como ausencia en el registro de asistencia cualquiera de los siguientes tres casos:

\begin{itemize}[leftmargin=*,itemsep=3pt,topsep=4pt]
    \item \textbf{Ingreso extemporáneo:} el caso de quienes lleguen quince minutos después de iniciada la sesión. En estos casos, el estudiante podrá permanecer y participar en la clase, pero no contará con un registro válido de asistencia.
    \item \textbf{Uso de cámara en modalidad PAT:} el caso de quienes, en modalidad PAT, no enciendan sus cámaras, salvo que se haya acordado previamente con el docente una excepción debidamente justificada.    
    \item \textbf{Ausencia de respuesta en espacios de participación:} el caso de quienes, una vez iniciada la sesión, sean llamados a participar y no respondan, salvo situaciones excepcionales.
\end{itemize}


% ------------------------------------------------------------
% COMUNICACIÓN
% ------------------------------------------------------------

\subsection{Estrategias de comunicación}

El programa, el cronograma detallado, las lecturas obligatorias y
los materiales complementarios estarán disponibles en la carpeta
compartida de Google Drive del curso.

La comunicación institucional oficial se realizará exclusivamente
mediante el correo electrónico
\href{mailto:\correoDocente}{\correoDocente}.

Durante la primera sesión se propondrá además la creación de un
grupo de comunicación rápida mediante WhatsApp, sujeto a la
aprobación de la mayoría del grupo.


% ------------------------------------------------------------
% HONESTIDAD ACADÉMICA
% ------------------------------------------------------------

\subsection{Honestidad académica y manejo de fuentes}

Todo trabajo escrito o entregable presentado en el marco de la
asignatura debe cumplir estrictamente con las normas de citación y
atribución de fuentes establecidas por APA, séptima edición.

Tanto en las intervenciones escritas como en las orales debe
mantenerse una distinción clara, consistente y veraz entre las
voces de los autores consultados y la voz propia.

Cualquier forma de plagio, entendida como la apropiación de textos
o ideas ajenas sin la debida acreditación, implicará una
calificación de cero punto cero (0.0) en la entrega correspondiente
y podrá ser remitida a las instancias disciplinarias de la
Coordinación del Programa, según la gravedad del caso.


% ------------------------------------------------------------
% INTELIGENCIA ARTIFICIAL
% ------------------------------------------------------------

\subsection{Uso crítico y transparente de inteligencia artificial}

En este curso no se prohíbe el uso de herramientas de inteligencia
artificial generativa, como ChatGPT, Claude, Gemini o similares.
Por el contrario, se incentiva su empleo como apoyo en los procesos
de aprendizaje y exploración, siempre que se realice desde una
postura analítica, consciente y éticamente responsable.

Su uso está sujeto a una regla de transparencia manifiesta: si un
taller, informe o presentación se apoyó en una herramienta de
inteligencia artificial, la entrega deberá incluir una
\textbf{Declaración de uso de IA} al final del documento.

% enumerate crea una lista numerada.
% label=\roman*) utiliza números romanos en minúscula.
\begin{enumerate}[
    label=\roman*),
    leftmargin=*,
    itemsep=3pt
]

\item
La plataforma y el modelo utilizados.

\item
Las tareas específicas para las que se empleó la herramienta, como
revisión de estilo, contraste de hipótesis, generación de
contraejemplos o estructuración de código o tablas.

\item
Los prompts o instrucciones centrales empleados, cuando resulte
pertinente.

\end{enumerate}

El uso irreflexivo, acrítico u oculto de inteligencia artificial,
presentando textos o análisis generados automáticamente como
producción propia sin la correspondiente declaración, se
considerará una falta a la honestidad académica y será evaluado con
la nota mínima de cero punto cero (0.0)
